\documentclass[12pt]{article}

\usepackage[a4paper, margin=1in]{geometry}
\usepackage{amsmath,amssymb,amsthm}
\usepackage{mathtools}
\usepackage{mathrsfs}
\usepackage{enumitem}
\usepackage{xcolor}
\usepackage[hidelinks]{hyperref}
\usepackage{fontspec}

\defaultfontfeatures{Ligatures=TeX}
\IfFontExistsTF{Times New Roman}
  {\setmainfont{Times New Roman}}
  {\setmainfont{TeX Gyre Termes}}
\IfFontExistsTF{Menlo}
  {\setmonofont{Menlo}}
  {\setmonofont{Latin Modern Mono}}

\newtheorem{theorem}{Theorem}[section]
\newtheorem{proposition}[theorem]{Proposition}
\newtheorem{lemma}[theorem]{Lemma}
\newtheorem{corollary}[theorem]{Corollary}
\theoremstyle{definition}
\newtheorem{definition}[theorem]{Definition}
\newtheorem{example}[theorem]{Example}
\theoremstyle{remark}
\newtheorem{remark}[theorem]{Remark}

\newcommand{\tightlist}{%
  \setlength{\itemsep}{0pt}\setlength{\parskip}{0pt}}

\title{%
  \textbf{Foundation II: Shared Manifestation}\\
  \textbf{Shared Structure, Public Calibration and Different Histories}%
}

\author{
  Sidong Liu, PhD \\
  \small iBioStratix Ltd \\
  \small \texttt{sidongliu@hotmail.com}
}

\date{Revised October 2026}

\hypersetup{pdftitle={Foundation II: Shared Manifestation. Shared Structure, Public Calibration and Different Histories},pdfauthor={Sidong Liu}}
\begin{document}
\emergencystretch=2em
\raggedbottom
\hbadness=5000
\vbadness=5000

\maketitle

\begin{abstract}
Several readers of one process can have different accessible algebras while
sharing a world. Their common structure is the intersection of their stable
hulls. It is the largest algebra available to all of them, inherits the common
modular flow and, with a supplied positive translation structure and a unique
vacuum, contains a proper half-sided modular inclusion whenever it is
non-scalar. No foreground overlap or commutation of noise channels is required.
For compact symmetry-selected hulls the common algebra is fixed by the group
generated by the readers' stabilizers; under the minimal-action Galois
conditions their jointly generated algebra is fixed by the intersection of
those stabilizers. The SU(2) level-one model realizes a shared Virasoro sector
between two different current sectors, even with disjoint foreground seeds.

Agreement does not certify the truth of a shared description. A common symmetry
change of instruments and preparations is indistinguishable, in every finite
adaptive record law, from the corresponding change of actual state. We derive
the effective groups of these errors and show that their kernels differ on
individual and shared content. Conversely, a non-scalar translation-stable
shared algebra contains a finite positive-energy pulse. With preparation and
readout permission, a known vacuum baseline turns it into an energy-orientation
reference with an explicit finite-record error bound. Under the additional
source regularity and modular-consistency conditions of Consciousness~II, the
pulse also supplies a spectral witness for a common sourced geometric response.
It does not fix a metric without boundary or initial data. For a finite
family of common frame errors, legal reference blocks identify a public target
exactly when indistinguishable reference laws have the same target value.
The resulting finite-record bound permits residual private frame differences;
an explicit reference protocol preserves different old histories. A
two-observer calibration model measures common error and fault transfer.
Independent-source compatibility is treated separately. Shared carrier,
shared state, agreeing records and correct content are distinct requirements.
\end{abstract}

\clearpage
\tableofcontents
\clearpage

\section{Introduction}\label{sec:intro}

Foundation~I asks how a source, its noise and a retained seed specify an
accessible algebra~\cite{FoundationI}. This paper asks what several readers
of that source share. The question has two parts. A common carrier must be
constructed from the readers' access. A common description must also survive
errors that all readers can make together. Agreement alone addresses neither
part completely.

The main setting is one ambient process with different seeds and accessible
hulls. Their intersection derives the shared algebra instead of assuming a
candidate carrier. Its invariance needs no commutation of the noise channels.
If a positive translation structure is supplied, a non-scalar intersection
inherits a Cut on its own vacuum-cyclic space. Cuts sit inside that shared
algebra; the algebra itself is not asserted to be a Cut of the ambient algebra.
The translations remain physical inputs, as in Foundation~I and
Consciousness~I~\cite{FoundationI,PaperI}.

Symmetry makes the distinction between common and joint access precise.
For hulls fixed by subgroups, the shared algebra is fixed by the group those
subgroups generate. The joint algebra obeys a different law under the compact
Galois hypotheses~\cite{ILP1998}. Two independent current seeds in the SU(2)
level-one net have no foreground overlap yet share a Virasoro sector. The
algebraic sharing mechanism is therefore stronger than foreground matching.
Foundation~III provides a noncommuting, repeated-noise realization of the same
fixed intersection~\cite{FoundationIII}; the present paper uses that result
without repeating its convergence proof.

Readers can then agree and still share a displaced standard. Reference-frame
restrictions and relational encodings are established subjects in quantum
information~\cite{BRS2007}. Here an adaptive-protocol covariance identity
identifies a class of common errors that no agreement statistic can expose.
The group acting faithfully on individual content can be larger than the group
acting faithfully on shared content. In the current-net model a common energy
sector is insensitive to reversals of current orientation. This is an
operational distinction, not a claim that a symmetry-invariant reading is
automatically true.

A shared algebra can also contain a way to test a standard. Consciousness~II's
bounded-preparation construction gives a finite positive-energy pulse whenever
the shared algebra is non-scalar and translation-stable~\cite{PaperII}.
If the pulse and vacuum baseline can actually be used at a common energy port,
the mean difference determines an otherwise unknown energy-gain sign.
The preparation changes the physical state, so it can also provide the source
witness required to identify a geometric response, subject to the separate
regularity and consistency conditions. This connects selected access to a
testable world rather than identifying geometry from consensus alone.

The source need not have been described by one state in every problem.
Independent source states, projections and memory kernels have a different
compatibility question. Appendices~\ref{sec:datum}--\ref{sec:nsource} retain
that theorem and its nonempty examples. It cannot be made unconditional by
substituting the one-process construction for its independent states.
Likewise, the two-observer simulation of Section~\ref{sec:model} is a
finite-resource calibration experiment, not a finite reconstruction of a
Type~III algebra.

Quantum Darwinism and spectrum broadcast structures study information held
redundantly in environmental fragments~\cite{Zurek2009,BrandaoPianiHorodecki2015,KorbiczHorodeckiHorodecki2014}.
An intersection of accessible algebras is not a redundancy certificate, and
two carriers can encode the same bit even when their literal algebraic
intersection is scalar. The paper keeps these distinctions while connecting
the source-selection problem of HAFF to the reference and self-calibration
problems of T-DOME~\cite{HAFFA,HAFFD,TDOMEII,TDOMEIII}.

Section~\ref{sec:onesource} derives shared structure and its symmetry laws.
Section~\ref{sec:commonerror} identifies common frame ambiguities.
Section~\ref{sec:sharedsource} constructs the internal reference and its
conditional geometric test. Section~\ref{sec:model} gives the calibration
simulation, and Section~\ref{sec:open} returns the results to the full theory.

\section{One process, several accessible structures}\label{sec:onesource}

Let $M$ be a von Neumann algebra with a cyclic and separating vector $\Omega$
and faithful normal reference state $\omega(x)=\langle\Omega,x\Omega\rangle$.
A family $\mathfrak E$ of normal unital completely positive maps preserves $\omega$.
Its members need not commute. Include the full reference modular group and put
\[
 \Gamma_2=\{\sigma_s^\omega:s\in\mathbb R\}\cup\mathfrak E.
\]
When a positive translation group $U(a)=e^{iaP}$ is supplied, with $P\ge0$,
$U(a)\Omega=\Omega$ and $U(a)MU(a)^*\subseteq M$ for $a\ge0$, also put
\[
 \Gamma_3=\Gamma_2\cup\{\operatorname{Ad}U(a):a\ge0\}.
\]
For either family, a subalgebra is admissible when every map in the family takes
it into itself. A reader retains a seed $S_j\subseteq M$ and has the ideal
accessible hull $B_j=\operatorname{Hull}_\Gamma(S_j)$: the intersection of all
admissible von Neumann subalgebras containing $S_j$~\cite{FoundationI}.
This closure describes an operator structure. It is not a finite-resource
reconstruction algorithm or a licence to implement every unitary in the hull.
The seed, noise, control permissions and reference state are physical inputs.

\begin{theorem}[The shared structure of one process]\label{thm:onesource}
For any family of readers with the same $M$, $\omega$ and $\Gamma$, set
$B_\cap=\bigcap_jB_j$.
\begin{enumerate}
\item $B_\cap$ is the largest admissible von Neumann subalgebra contained in
every $B_j$. It contains $\operatorname{Hull}_\Gamma(T)$ for every
$T\subseteq B_\cap$, and carries a unique normal $\omega$-preserving conditional
expectation. Its modular group in $\overline{B_\cap\Omega}$ is
$\sigma^\omega|_{B_\cap}$.
\item For every normal actual state $\psi$ of $M$, the restrictions
$\psi|_{B_j}$ agree on $B_\cap$. This says nothing about equality of the readers'
joint records or about the accuracy of their decoding.
\item For $\Gamma_3$, suppose $U$ is the translation group of an ambient
half-sided modular inclusion and $\ker P=\mathbb C\Omega$. If
$B_\cap\ne\mathbb C1$, then $U(a)B_\cap U(a)^*\subset B_\cap$ for $a>0$
is a proper half-sided modular inclusion on $\overline{B_\cap\Omega}$.
The recovered translations are the restricted ambient ones. If the inclusion
is standard, Consciousness~I's Cut--geometry correspondence applies there.
\end{enumerate}
\end{theorem}
\begin{proof}
An intersection of von Neumann algebras is a von Neumann algebra. If $x$ is
in every $B_j$, then $\gamma(x)$ is in every $B_j$ for each $\gamma\in\Gamma$.
For the modular maps, including both $s$ and $-s$ gives equality rather than
only inclusion. The intersection is therefore admissible, and every admissible
subalgebra contained in all the $B_j$ lies in it. If $T\subseteq B_\cap$, the
intersection is itself a competitor in the definition of its hull.
Modular invariance gives the expectation and restricted modular group by
Takesaki's theorem~\cite{Takesaki1972,Hiai2021}. Restrictions of one state agree
on an intersection. Finally $B_\cap$ inherits forward translation invariance;
Foundation~I's Cut-inheritance theorem supplies (3), including properness from
the unique invariant vector hypothesis~\cite[section on Cuts in the hull]{FoundationI}.
\end{proof}

This route needs neither a candidate shared carrier nor commuting projected
memory kernels. It also does not require the foreground seeds themselves to
overlap. It addresses readers of one process. Independently specified source
states can still disagree on an overlap; their compatibility is treated in
Appendices~\ref{sec:datum}--\ref{sec:cp}. The reference $\omega$ determines the
closure and the Cut. The actual state $\psi$ may be different, even nonstationary.
Its restrictions agree because they are restrictions of one state, not because
the reference and actual states have been identified.

\subsection{Shared, individual and jointly available algebras}

Let a compact group $G$ act continuously on $M$, preserving $\omega$, and let
$K_j$ be closed subgroups such that $B_j=M^{K_j}$. Write
\[
 K_\cup=\overline{\langle\bigcup_jK_j\rangle},\qquad
 K_\cap=\bigcap_jK_j.
\]
The fixed-algebra identity
\begin{equation}\label{eq:sharedgroup}
 B_\cap=M^{K_\cup}
\end{equation}
holds without a minimal-action assumption: an element fixed by every $K_j$ is
fixed by their generated group and, by continuity, its closure.
If the action is minimal on a factor with separable predual, the compact-group
Galois correspondence also gives
\begin{equation}\label{eq:jointgroup}
 \bigvee_jB_j=M^{K_\cap}.
\end{equation}
Here minimal means faithful with $(M^G)'\cap M=\mathbb C1$~\cite{ILP1998}.
Indeed the join is an intermediate algebra $M^L$. Containment of every $M^{K_j}$
forces $L\subseteq K_\cap$, whereas the join lies in $M^{K_\cap}$, forcing the
opposite inclusion. If $M^{K_\cap}$ is also $\Gamma$-admissible, the hull of the
joint seed equals that join. Admissibility of a join is not automatic for
arbitrary UCP maps, which need not preserve products of elements from different
subalgebras.

\begin{example}[Shared geometry without shared foreground]\label{ex:twocurrents}
In the SU(2) level-one current net let $M=\mathcal A(0,\infty)$, with its
vacuum, modular dilations and forward translations. With no additional noise,
the bounded current seeds $\frac2\pi\arctan J_{e_3}(f)$ and
$\frac2\pi\arctan J_{e_1}(f)$ have hulls $M^{H_z}$ and $M^{H_x}$, where $H_z$
and $H_x$ are the axial rotation groups. These seed-to-hull statements and the
fixed subnet are established in Foundation~I~\cite{FoundationI}.
Here $f$ is nonzero, real and smooth with compact support in the positive
half-line. The two groups generate SO(3), so
\[
 B_\cap=M^{SO(3)}=\mathrm{Vir}_1(0,\infty),\qquad
 B_1\vee B_2=M.
\]
Each individual hull is a current algebra; their common algebra is Virasoro;
their joint algebra is the full current net. They are three different objects.
The two seed lines have zero foreground intersection, yet their hulls share
a standard Cut. Foreground overlap is therefore not necessary for sharing.
The supplied vacuum net and translations remain inputs, not outputs of the
sharing argument.
\end{example}

Nontriviality is not universal. In $M_2\otimes M_2$ with the trace, trivial
modular flow and no noise, the seed algebras $W^*(Z\otimes1)$ and
$W^*(1\otimes Z)$ intersect only in the scalars. The two registers can nevertheless
encode the same classical bit with high agreement. A shared subalgebra on one
ambient carrier and matching content on different carriers are distinct ways
of sharing; neither defines the other.

\subsection{A dynamical realization without commuting noises}

For a finite family $K_1,\ldots,K_m$ let $Q_j$ be the Haar expectations and apply
$D_n=(1-\theta_n)\mathrm{id}+\theta_nQ_{i_n}$, with $0<\theta_n\le1$.
Foundation~III proves
\begin{equation}\label{eq:shareddynamics}
 \sum_{n:i_n=j}\theta_n=\infty\quad\text{for each }j
 \quad\Longrightarrow\quad
 D_n\cdots D_1(x)\longrightarrow Q_{K_\cup}(x)\quad\text{strong*}.
\end{equation}
Thus $B_\cap$ can be the algebra retained by repeated noise, not only an
intersection written after the fact~\cite{FoundationIII}.
There is no bounded-gap or commutation condition. Record-dependent schedules
are covered when their realized sequence meets the exposure condition, but
intervening measurement and control channels must be included separately.
The finite-family restriction is essential; Foundation~III supplies a countable
family for which every noise is visited indefinitely but convergence fails.
This constructs persistence for supplied noise actions. It does not derive why
the source has those actions, or turn the ideal limit into a finite-time
resource guarantee.

\section{Common errors that agreement cannot expose}\label{sec:commonerror}

A reader's standard assigns record labels to instruments, preparations and
decoders. Correlating standards is a reference-frame problem, not the equality
of two restricted states~\cite{BRS2007}. Let $\alpha_g$ be a normal automorphism
of $M$ commuting with every uncontrolled evolution map used in a protocol.
Each instrument branch $\mathcal I^y$ is a normal completely positive map,
and its sum over outcomes is unital. Choices of later branches and preparations
may depend on earlier labels and on recorded private randomness, with the same
prescribed conditional coin law in the two configurations being compared.

\begin{proposition}[An indistinguishable common change of standard]
\label{prop:recordcovariance}
For any finite adaptive protocol, replacing every controlled operation by
$\alpha_g\mathcal I^y\alpha_g^{-1}$ in the actual state $\psi$ gives the same
joint record law as the original protocol in $\psi\circ\alpha_g$.
If $x=x^*\in M$, $\alpha_g(x)=-x$ and $\psi(x)\ne0$, a binary report about the
sign of $\psi(x)$ has correctness probabilities adding to one in these two
indistinguishable configurations. At least one of them has correctness at most
$1/2$, regardless of the agreement rate or the number of readers.
\end{proposition}
\begin{proof}
For a fixed record string the probability is a state applied to the ordered
composition of instrument branches and intervening evolution maps, evaluated
at $1$. Conjugating the controlled branches and commuting $\alpha_g$ through
the uncontrolled maps leaves $\alpha_g$ on the outside and
$\alpha_g^{-1}(1)=1$ on the inside. This changes $\psi$ to
$\psi\circ\alpha_g$. Adaptivity changes the branch selected by the fixed
string, not this identity. Averaging over private coins preserves it.
The truth changes sign between the two configurations while the report law
does not. A binary report is correct in exactly one of them, proving the sum
and the bound.
\end{proof}

The ambiguity concerns two physical configurations: a displaced standard and
a displaced state. It is not the assertion that the same state has two
different true values. A calibrated reference left unchanged by the protocol
can break the symmetry; conjugating that reference as well would remove the
very information that makes it a reference. This result classifies a source
of common error, not every statistical mistake. Two independent error bits
with error probability $e$ still agree wrongly with probability $e^2$.

\subsection{The effective group of a common frame error}

Let $G$ be a compact group of allowed frame transformations, and let $H\le G$
be its closed subgroup preserving the uncontrolled world laws.
Suppose the individual algebras are $B_j=M^{K_j}$. We first consider errors
that preserve each labelled reader's algebra rather than exchange the readers.
Define
\[
 N=H\cap\bigcap_j N_G(K_j),
 \qquad L=N\cap\bigcap_jK_j,
 \qquad L_\cap=N\cap K_\cup .
\]
Under the minimal compact-action hypotheses used in~\eqref{eq:jointgroup}, the
pointwise stabilizer of $M^K$ inside $G$ is exactly $K$~\cite[Remark~3.17]{ILP1998}.

\begin{proposition}[Joint and shared errors have different kernels]
\label{prop:errorgroup}
In this regime the transformations in $H$ preserving every $B_j$ are exactly
$N$. Their faithful action on the tuple of individual algebras, or on their
jointly generated algebra, is $N/L$. Their faithful action on the shared algebra
$B_\cap=M^{K_\cup}$ is $N/L_\cap$. In particular $L\subseteq L_\cap$, so a
nontrivial error on individual content can be invisible on all shared content.
For one reader the effective group is
\[
 N_H(K)/(K\cap H),
\]
where $N_H(K):=H\cap N_G(K)$ even when $K$ is not contained in $H$.
This reduces to $N_H(K)/K$ only when $K\subseteq H$.
For one reader the observables unchanged by every such error form
$M^{\overline{\langle K,N_H(K)\rangle}}$.
\end{proposition}
\begin{proof}
$\alpha_g(M^K)=M^{gKg^{-1}}$, so injectivity of the fixed-subgroup correspondence
makes preservation of $M^K$ equivalent to normalization of $K$.
The pointwise kernel on each $B_j$ is $N\cap K_j$; intersecting gives $L$.
Fixing every individual algebra pointwise is equivalent to fixing its join
pointwise. Every element of $N$ also normalizes $K_\cup$, and the pointwise
kernel on $M^{K_\cup}$ is $N\cap K_\cup=L_\cap$.
Both are kernels of continuous group actions, hence closed normal subgroups
of $N$, even if $K_j$ is not contained in $H$. Also $K_\cap\subseteq K_\cup$.
The invariant-algebra statement follows by imposing invariance under $K$ and
$N_H(K)$ together.
\end{proof}

Without the Galois hypothesis one can still restrict normalizing transformations
to the reader's algebra, but the effective group is the group divided by its
actual pointwise kernel. That kernel need not be $K\cap H$.
The invariant algebra in Proposition~\ref{prop:errorgroup} is immune to the
specified frame errors. Its readings can still have statistical error or be
decoded incorrectly; immunity is not correctness.

\begin{example}[Energy is shared while a current can be reversed]
In the single-axis SU(2) level-one noise model the world-law symmetries inside
SO(3) form the axial normalizer $O(2)_z$, while a current reader has
$K=H_z\cong SO(2)$. The residual frame group is $\mathbb Z_2$: its nontrivial
element reverses $J_z$ but fixes the stress tensor. The immune algebra is the
$O(2)_z$ orbifold, not the full current algebra~\cite{FoundationI}.
For the two perpendicular current readers of Example~\ref{ex:twocurrents},
with no extra noise, $N$ is the four-element group of identity and the three
half-turns about the coordinate axes. Its action on the pair of private
algebras is faithful, whereas it acts trivially on their shared Virasoro
algebra. The energy sector is unchanged while either current orientation can
be wrong. This gives an explicit nontrivial $N/L$ with trivial $N/L_\cap$.
\end{example}

\section{A reference made inside the shared structure}\label{sec:sharedsource}

Agreement cannot break a symmetry of the whole protocol. A preparation with a
known baseline can. The distinction matters for the connection from algebra
selection to geometry: a changed reading is not a changed material source,
whereas an implemented excitation is.

\begin{proposition}[Shared positive-energy preparation and orientation]
\label{prop:sharedreference}
Suppose the readers' hulls have the common representation and $\Gamma_3$
translation closure of Theorem~\ref{thm:onesource}, with
$\ker P=\mathbb C\Omega$. If $B_\cap$ is non-scalar, it contains a bounded
self-adjoint non-scalar operator $B$ obtained by smoothing a translation orbit,
and a unitary $W_\theta=e^{i\theta B}$ with
\[
 0<E_\theta=\langle W_\theta\Omega,PW_\theta\Omega\rangle<\infty
 \quad\text{for }0<|\theta|\|B\|<\pi.
\]
All energy moments of this state are finite. If this unitary is an allowed
preparation and the same energy port can be used on the known vacuum and on
freshly prepared states, suppose its output is
\[
 Y=\epsilon g\,p+b+\nu,\qquad \epsilon\in\{-1,1\},\quad g>0,
\]
where $p$ is the spectral energy outcome and the noise has the same finite
mean in both uses. The difference of population means is
$\epsilon gE_\theta$, so these record laws identify the sign $\epsilon$.
\end{proposition}
\begin{proof}
The intersection inherits separation and the forward translation orbit of
every one of its operators. Consciousness~II's bounded-preparation proposition
therefore applies~\cite[section on a bounded accessible preparation]{PaperII}:
for a narrow real nonnegative $f\in C_c^\infty(0,\tau)$ with $\int f=1$ and
a bounded non-scalar $C=C^*\in B_\cap$, set
$B=\int f(t)U(t)CU(t)^*dt$. As the support of $f$ shrinks to $t_0$,
$B\Omega\to U(t_0)C\Omega\notin\mathbb C\Omega$ by separation,
so a sufficiently narrow $f$ keeps $B$ non-scalar. Its bounded iterated
commutators with $P$ then give
finite energy moments of $W_\theta\Omega$. Separation and the unique ground
vector give strictly positive energy for non-scalar $W_\theta$.
The vacuum has $p=0$, so subtracting its output mean removes both the unknown
offset and the common noise mean. The remaining sign is $\epsilon$.
\end{proof}

$P$ is the common translation generator of the model. In a null Cut it is
null momentum, not automatically the time-translation Hamiltonian of a
higher-dimensional source. Positivity alone does not supply this calibration:
in a finite two-level system $P=\operatorname{diag}(0,1)$ satisfies $XPX=1-P$.
The reversal fails to preserve the known vacuum baseline, which the preceding
protocol requires. Reading the baseline and implementing the excitation are
real permissions, not consequences of algebra membership.

\begin{corollary}[A finite-record orientation certificate]
\label{cor:orientationcost}
Assume independent fresh samples and additive noise independent of the energy
outcome, with variance $\sigma_\nu^2$ common to the two uses. From $n$ vacuum
reads and $n$ prepared reads, the sign of their sample-mean difference has
error probability at most
\begin{equation}\label{eq:orientationcost}
 \min\left\{1,\frac{g^2\operatorname{Var}_\theta(P)+2\sigma_\nu^2}
 {n g^2E_\theta^2}\right\}.
\end{equation}
This protocol uses $2n$ reads and $n$ fresh preparations, whose prepared
subsystem energy is $nE_\theta$. Switching, detection, reference preparation
and reset costs are separate. A uniform prescribed-confidence budget needs
bounds on the gain, noise and energy variance.
\end{corollary}
\begin{proof}
The mean difference has mean $\epsilon gE_\theta$ and variance
$(g^2\operatorname{Var}_\theta(P)+2\sigma_\nu^2)/n$. A wrong sign requires a
deviation at least $gE_\theta$; Chebyshev's inequality gives the bound.
\end{proof}

This reference breaks energy-orientation ambiguity but not every internal
frame symmetry. A group commuting with $P$ can still reverse a charge or
current. Thus a common positive pulse calibrates one part of the world without
supplying a universal external standard.

\subsection{Public correctness without equal private histories}
\label{sec:publictarget}

Calibrating a public target need not identify the whole frame. Let
\(F\) be a finite set of possible common frame errors. A legal reference
block comprises specified preparations, their source changes, instrument
uses and recorded coins. It has a finite record alphabet and known law
\(p_f(z)\) under \(f\in F\). These laws concern actual reference operations
that are held fixed when comparing frames. They are not inferred from the
agreement of the readers. Let \(t(f)\) be a finite-valued public target.

\begin{proposition}[Target-specific reference criterion]
\label{prop:targetcalibration}
Repeated independent fresh reference blocks permit uniformly consistent
estimation of \(t(f)\) if and only if
\[
 p_f=p_{f'}\ \Longrightarrow\ t(f)=t(f').
\]
When the criterion holds, put
\[
 r=\max_{t(f)\ne t(f')}
       \sum_z\sqrt{p_f(z)p_{f'}(z)}<1.
\]
For \(n\) blocks, a maximum-likelihood estimate \(\hat f\), followed
by the report \(t(\hat f)\), has target error at most
\[
 \min\{1,(|F|-1)r^n\}.
\]
If no two hypotheses have different targets the error is zero. A single
decoded public report transmitted without corruption makes every reader
correct on the same event, irrespective of their private histories.
\end{proposition}
\begin{proof}
If two identical laws have different targets, their complete \(n\)-block
laws are identical for every \(n\). The probabilities of reporting their
two different correct values sum to at most one, so one error is at least
\(1/2\). Conversely, under true \(f\), selection of a wrong-target
candidate \(f'\) requires its likelihood to be at least that of \(f\).
Markov's inequality applied to the square root likelihood ratio bounds
this event by
\(\bigl(\sum_z\sqrt{p_f(z)p_{f'}(z)}\bigr)^n\).
Zero-probability records are treated on the support of \(p_f\).
The union bound over wrong-target candidates gives the result.
Distinct probability laws have affinity below one by Cauchy--Schwarz.
Finiteness makes the maximum strictly below one.
\end{proof}

The testing inequality is standard; the structural consequence here is
that public calibration is determined by the target's indistinguishability
classes. A full frame reconstruction would require separation of every
distinct frame. When frame symmetries are represented by a group,
unremoved transformations may still act on private currents while fixing
the public energy target. Identifying a reference law is also weaker than
implementing the inverse frame transformation: reporting and control are
separate permissions.

There is a concrete finite physical witness. Let a reference qubit have
\(H_R=\Delta|1\rangle\langle1|\), \(\Delta>0\). Legal preparations
\(|0\rangle,|1\rangle\) and an energy pointer supply the known baseline
and excited input. Suppose the pointer's two values are reversed by an
unknown common sign \(f=\pm1\) and its binary read flips with probability
\(e<1/2\). Reading the excited input gives Bernoulli laws with
affinity \(2\sqrt{e(1-e)}\), so \(n\) reads identify the common sign
with error at most \([2\sqrt{e(1-e)}]^n\).
An ideal controlled pointer interaction realizes the unflipped
measurement; the read noise is a supplied physical channel.
The source preparation uses \(n\Delta\) in excited-subsystem energy.
The pointer preparation, coupling, elapsed time, storage and disposal of
spent systems have their own costs. Energy supplied is not net heat.

Let \(D\) contain all old private memory registers and their old
references, in an arbitrary state \(\rho_D\), with different histories
at different readers. Prepare the reference blocks independently of
\(D\) and apply the calibration instrument only to those blocks.
The final state is exactly
\(\rho_D\otimes\rho_{\rm calibration}\); all old correlations in
\(D\) remain intact, even conditional on the calibration outcome.
More generally, for initially correlated active and old registers,
a trace-preserving map acting only on the active register preserves the
old marginal nonselectively, whereas a selected new record can reweight
the old history. This is the distinction needed by the retained-history
return in Foundation~III \cite{FoundationIII}.
The tensor architecture is an existence witness, rather than a
factorization theorem for Type~III algebras.

The reference criterion fails if the excitation and its label are
conjugated along with the faulty instruments: Proposition~\ref{prop:recordcovariance}
then applies again. It also fails for a destroyed pointer with \(e=1/2\).
A shared algebra supplies possible operators; it supplies neither an
independently trusted reference nor the preparation and communication
operations used here. For a reference block costing resource vector
\(\mathbf c\), \(n\) blocks cost \(n\mathbf c\), in addition to
the retained private storage and public communication. Those requirements
make the claimed public correctness falsifiable.

\subsection{From a shared preparation to a geometric test}

Suppose an internal compact group $G$ fixes the bounded functions of the
required stress-tensor smearings. Every selected algebra of the form $M^{K_j}$,
$K_j\le G$, then contains those readings, as does $B_\cap$. This is a direct
condition for a common stress sector; it is not a theorem about every non-scalar
hull. Readers of separate decoupled fields can have disjoint stress sectors.

\begin{proposition}[A shared test of the sourced geometric response]
\label{prop:geometrytest}
In Consciousness~II's conformal Minkowski setting, suppose a shared
positive-energy pulse is implemented, either the bounded pulse of
Proposition~\ref{prop:sharedreference} or an explicit scalar preparation
of that paper. Its translation generator
is the physical positive Hamiltonian, and its vacuum-subtracted source has an
integrable energy density with integral $E_\theta>0$. Suppose the source and
response belong to the smooth tempered Radon Fourier class, and the full
centre, open-frame and modular-consistency hypotheses of that paper hold.
If every reader has the required stress and boundary observables, uses the same
actual state and normalization, and applies the same geometric readout law,
the pulse has spectral weight away from zero and identifies their common
canonical response up to the law's normalization. They obtain the same
matter-sourced linearized Einstein tensor. A unique metric requires additional
common initial or boundary data and gauge choices.
\end{proposition}
\begin{proof}
If a Radon Fourier measure were supported only at zero, the source would be
constant. Integrability of its energy density on the infinite spatial slice
would force that density to vanish, contrary to $E_\theta>0$.
The spectral identification theorem and common sourced response of
Consciousness~II now apply~\cite{PaperII}. A homogeneous gravitational wave in
dimension at least four can have zero linearized Einstein tensor and nonzero
linearized curvature, so the source alone cannot fix the metric.
\end{proof}

Finite energy does not itself imply the Radon regularity used here, and the
finite-amplitude pulse need not have a nonzero first-order stress tangent at
zero amplitude. Nor does the chiral SU(2) model establish all the
higher-dimensional hypotheses. The bridge states exactly which source and
readout data must be common. A unitary used for calibration physically changes
the state; merely choosing a different label or algebra does not. The
implemented source and its costs are the link to HAFF-D~\cite{HAFFD}; bounded
source feedback and persistent access are treated in
Foundation~III~\cite{FoundationIII}, rather than inferred from agreement.

\begin{example}[Shared hulls and a physical time generator]
\label{ex:sharedphysicalH}
Take two independent massless conformal complex scalar fields in
$3+1$ dimensions, on one vacuum Fock space, and let
$M=\mathcal A(V_+)=\mathcal A_1(V_+)\bar\otimes\mathcal A_2(V_+)$.
The future-cone vacuum algebra is standard and its modular group is
dilation \cite[Theorem~2.3(ii)]{BGL1993}; the free massless-field
modular structure was studied in \cite{HislopLongo1982}.
Write $\mathcal N_j(O)=\mathcal A_j(O)^{U(1)}$ and give the two readers
the selected nets
$B_1(O)=\mathcal N_1(O)\bar\otimes\mathcal A_2(O)$,
$B_2(O)=\mathcal A_1(O)\bar\otimes\mathcal N_2(O)$.
For $O=V_+$ their shared algebra is
$\mathcal N_1(V_+)\bar\otimes\mathcal N_2(V_+)$.
The two vacuum-preserving Haar expectations commute, proving this
intersection directly. Specify the noise family to be
$\mathcal E=\{Q_1\otimes\mathrm{id},\mathrm{id}\otimes Q_2\}$,
where $Q_j$ is Haar projection on field $j$. Both maps preserve
each $B_i(V_+)$: one acts identically on its fixed factor, and the
other projects its unrestricted factor into a subalgebra.
The bounded self-adjoint parts of these algebras can serve as ideal
seed families. Their invariance under this noise family, the full
modular group and forward time translations makes their $\Gamma_3$
hulls exactly the displayed cone algebras.

Here $U(t)=e^{it(H_1+H_2)}$ is the physical time translation,
$V_++te_0\subset V_+$ for $t\ge0$, and its only zero-energy vector
is the vacuum. The latter follows from strictly positive
one-particle energy $|p|$ almost everywhere. Gauge-fixed algebras
are represented on their own vacuum-cyclic subspaces, not asserted
cyclic on the ambient Fock space. Haar projection commutes with $H$;
its restriction to each such subspace retains positivity, the same
physical translation and the unique vacuum.

A compact neutral quadratic preparation in field~1, as specified in
Consciousness~II, lies in both readers' algebras, has strictly
positive finite $H$ energy and a smooth nonzero-frequency Radon
source \cite[Proposition on a neutral preparation]{PaperII}.
Their selected algebras at every diamond contain the bounded stress
and collar readings. Haar expectations preserve the common vacuum;
permitted preparations, calibrated collar limits and common
$\mathrm{MC}_\eta$ normalization remain inputs. The constraint-compatible
linear response of that paper then instantiates
Proposition~\ref{prop:geometrytest}.
The reference hull construction uses the cone algebra; the geometric
test uses the same gauge-fixed nets on all $D(B)$. This supplies the
required diamond readings, not a claim that each diamond algebra is
the cone's $\Gamma_3$ hull or that every ball operator belongs to
$\mathcal A(V_+)$. The generic bounded pulse also
exists here, but its source regularity is not inferred solely from
its energy moments.
\end{example}
\section{Agreement, calibration and fault transfer}\label{sec:model}

A finite two-observer experiment tests the distinct roles of agreement
and paid references. A bank of $4096$ bits evolves under a hidden event
$c_k=\pm1$, flipping with probability $.01$ per epoch. Each bit is
retained with probability $.65$, replaced by a fair coin with probability
$.25$, or drawn with mean $.6c_k$ with probability $.10$.
Each drifting binary instrument has mean response $g_i z+b_i$ and a
$512$-read budget per epoch. Known $\pm1$ references displace world
reads; agreement exchanges current estimates and changes next-epoch
bias calibration. Only the scorer accesses $c_k$.
Appendix~\ref{app:protocol} and the supplementary protocol give the
legal clipping, seeds, updates and full outputs.

For binary reports put $A=P(\hat c_1=\hat c_2)$,
$C_i=P(\hat c_i=c)$ and $E=P(\hat c_1=\hat c_2\ne c)$.
Eventwise counting, without independence, gives
\[
 C_1+C_2=1+A-2E,\qquad
 P(\mathrm{correct}\mid\mathrm{agree})=1-E/A\quad(A>0).
\]
\begin{table}[htbp]
\centering
\begin{tabular}{lcccc}
\hline
Scheme & $A$ & $C_1$ & $E$ & $1-E/A$\\\hline
Initial references only & .647 & .753 & .067 & .896\\
Agreement only & .837 & .789 & .125 & .851\\
64 references per sign, no coupling & .860 & .908 & .020 & .977\\\hline
\end{tabular}
\caption{Selected comparisons under the same $512$-read budget:
$800$ trajectories of $3000$ epochs, excluding the first $200$;
correctness standard errors below $.002$. The supplementary protocol
retains all tested budgets and fault controls.}
\label{tab:agree}
\end{table}

Agreement rises while the reliability of an agreeing pair falls.
Among tested reference counts $4,16,32,64,128$ per sign, $64$ gives
the highest correctness; this is not a global optimum.
The full common-bias scan and unequal-reference controls are retained
with the protocol.

A $100$-epoch fault starts at a random epoch and adds a clipped $.30$
bias or sets one instrument's gain to zero. With $16$ references per
sign, unweighted agreement transfers zero-gain correctness losses
$.121,.123$ to the undamaged observer. Weighting by the partner's
reported gain reduces them to $.014,.016$; $64$ references reduce
them to $.003,.004$. A matched no-coupling observer has zero transfer.
A separately written implementation gives equal-weight harm $.110$
and weighted harm $.0036$. Isolation at a scorer-known fault time is
a diagnostic, not a learned deployment policy. References remove bias;
a zero-gain reading requires an additional gain-sensitive decision.

Communication costs $22$ bits per epoch for agreement and $44$ for
weighted agreement. All observers start with $4096$ reference reads.
The world is uncontrolled, the rules are designed, and only two readers
and the stated faults are tested. These data support calibration and
fault transfer in this protocol, not a consensus-to-geometric law.

\section{What makes a shared world testable}\label{sec:open}

The intersection fixes common observables, restrictions of one actual
state fix common state values, and instruments fix joint record laws.
Correct public reports additionally require a reference that separates
the public target's indistinguishability classes. This can leave
private frames and old histories different. The algebraic intersection
and the public-reference criterion answer complementary questions.

The common-pulse construction joins selected access to an actual
source test. Its physical translation, preparation, source regularity,
stress and collar readouts, and geometric consistency must coexist.
Example~\ref{ex:sharedphysicalH} supplies that algebra/source combination
in a conformal matter model. It retains the supplied coupling and
gravitational data; the null generator in the chiral example is a
different physical quantity. A sourced Einstein tensor does not
determine free radiation, a metric gauge, or experience.

We propose that the operational extent of a shared world is selected
partly by the references its participants can jointly sustain.
Source changes alter reference availability and price; changed
reference risk alters the legal common task family; the resulting
actions change the next source. Some common capabilities can fail
before private stored histories disappear. The finite reference
witness establishes calibration and old-history preservation.
Foundation~III supplies a finite semiclassical mean-action return.
One maintained reference/source/readout device remains a further
requirement, rather than a consequence of a strict hull inclusion.

Failure conditions are concrete. A shared hull can be scalar, a
reference can be conjugated with the faulty instruments, and a
zero-gain instrument can erase the information needed to correct it.
An energy reference cannot remove internal errors commuting with
energy. Paid controls, output entropy and reusable memory determine
the resource price; algebra size alone does not. Foundation~III,
T-DOME and Brain Anchor address the distinct return, calibration and
formation mechanisms. Independent-source state equality and a common
projected memory kernel retain their explicit appendix conditions.

\appendix

\section{Independent sources: compatibility and its limits}\label{sec:datum}

The same-process construction uses restrictions of one actual state. For
independent sources let \(M\) have a cyclic separating reference vector
\(\Omega\), reference state \(\omega_0\), faithful normal states
\(\varpi_i\), and normal UCP semigroups \(\mathcal E_t^{(i)}\).
Suppose each source supplies an individual algebra
\(A_i\subseteq M\), state \(\omega_i=\varpi_i|_{A_i}\), and realized
foreground space \(V_i\subseteq M_{\rm sa}\).
The independently chosen \(\varpi_i\) need not agree on \(A_1\cap A_2\).

For a candidate \(C\subseteq A_1\cap A_2\), the independent-source
conditions are:
\begin{description}
\item[JS-A.] \(C\) is invariant under both semigroups and under the full
reference modular group. The semigroups need not commute.
\item[JS-B.] A faithful normal state \(\psi\) on \(C\) is stationary
under both restricted semigroups and satisfies \(\omega_1|_C=\omega_2|_C=\psi\).
Equality and stationarity are substantive inputs.
\item[JS-C.] When a joint projected record process is needed, its normal
UCP idempotents \(P_1,P_2\) commute; their product has von Neumann-algebra
range containing \(C\), fixes \(C\), and the projected dynamics has one
specified common predual memory kernel. This last clause is a dynamical
condition, rather than a consequence of commuting projections
\cite{Nakajima1958,Zwanzig1960}.
\end{description}

\begin{proposition}[Commuting expectations]\label{prop:commuting_square}
If \(E_i:M\to M_i\) are commuting normal conditional expectations,
\(E_1E_2\) is a normal conditional expectation onto \(M_1\cap M_2\).
It supplies the algebraic clauses of JS-C when \(C\subseteq M_1\cap M_2\).
\end{proposition}
\begin{proof}
The product is normal UCP and idempotent. Its range lies in both
\(M_i\), and every element of their intersection is fixed by it. The
norm-one projection onto that subalgebra is a conditional expectation
\cite{AccardiCecchini1982,Popa1983}. No kernel identity follows.
\end{proof}

\section{An independent-source hull theorem}\label{sec:two}

The state equality is an input in JS-B. The following conclusion
constructs the minimal compatible hull inside a supplied candidate
$C$; it does not derive agreement of independent states.

For a nonzero finite-dimensional \(V_1\cap V_2\), choose a basis
\(O_1,\ldots,O_k\) in \(C\), and set \(S=W^*(1,O_1,\ldots,O_k)\).
The seed-space dimension is \(k=\dim(V_1\cap V_2)\le
\min(\dim V_1,\dim V_2)\). This bounds neither the algebraic dimension
of \(S\) nor the size of its hull. A single observable
\(O=\operatorname{diag}(0,1,2)\) in \(M_3\) spans a one-dimensional
seed space but generates the three-dimensional diagonal algebra.
One observable with infinite spectrum can likewise generate an
infinite-dimensional algebra. The finite seed bound also says nothing
about nontriviality. Let
\[
 \Phi=\bigcap\{N\subseteq M:N\text{ is a von Neumann algebra},
 S\subseteq N,\ \sigma_s^{\omega_0}(N)=N,\ 
 \mathcal E_t^{(i)}(N)\subseteq N\}.
\]

\begin{theorem}[Independent-source compatibility]\label{thm:main}
If \(C\) satisfies JS-A and JS-B and contains \(S\), then
\[
 S\subseteq\Phi\subseteq C\subseteq A_1\cap A_2,
 \qquad \omega_1|_\Phi=\omega_2|_\Phi=\psi|_\Phi .
\]
The state is faithful, and \(\Phi\) inherits the reference modular
flow on \(\overline{\Phi\Omega}\). If a proper half-sided modular
inclusion \(N\subset\Phi\) is supplied there, with \(\Omega\) cyclic
and separating for \(N\), Consciousness~I's Cut correspondence applies.
Standardness gives the net correspondence. A joint projected record
process additionally requires JS-C.
\end{theorem}
\begin{proof}
JS-A makes \(C\) a competitor in the defining intersection, so the
class is nonempty and \(\Phi\subseteq C\). Intersections preserve the
algebra, the seed inclusion and all specified invariances, proving
minimality. Restriction preserves faithfulness, and JS-B gives equality
of the states. Takesaki's theorem identifies the restricted reference
modular group \cite{Takesaki1972,Hiai2021}. The Cut and the projected
kernel are the separate supplied inputs just stated.
\end{proof}

The Cut sits inside \(\Phi\) in its reference representation; the hull
is not asserted to be a Cut of the ambient algebra. The reference
state \(\omega_0|_\Phi\) need not equal \(\psi|_\Phi\). The theorem
does not turn agreement into an output when JS-B already assumes it.

\begin{proposition}[Finite compatibility witness]\label{prop:fdwitness}
JS-A/B/C can coexist with disagreement elsewhere in the overlap.
\end{proposition}
\begin{proof}
In \(M_2(S)\otimes M_2(P_1)\otimes M_2(P_2)\) with tracial reference,
let \(D_X(a)=(a+Z_XaZ_X)/2\), \(P_1=D_SD_{P_1}\),
\(P_2=D_SD_{P_2}\), and
\(\mathcal E_t^{(i)}=e^{-\lambda_i t}{\rm id}
 +(1-e^{-\lambda_i t})P_i\), \(\lambda_i>0\).
Choose \(C=W^*(Z_S)\), \(A_1=A_2=W^*(Z_S,Z_{P_1},Z_{P_2})\), and
\[
 \varpi_1=\rho(1/3)\otimes\rho(1/4)\otimes\rho(1/5),\quad
 \varpi_2=\rho(1/3)\otimes\rho(-1/4)\otimes\rho(-1/5),
 \qquad \rho(a)=(1+aZ)/2 .
\]
All states are faithful. Both semigroups fix \(C\) pointwise, the
reference flow is trivial, and the states agree on \(C\), verifying
JS-A/B. The projections commute and their product has diagonal range.
On the fixed sector the common kernel is zero, verifying JS-C.
The foregrounds \(\operatorname{span}(Z_S,Z_{P_1})\) and
\(\operatorname{span}(Z_S,Z_{P_2})\) overlap in \(Z_S\), while the
state discrepancy on \(Z_{P_1}\) is \(1/2\). This finite witness does
not realize a Type~III Cut.
\end{proof}

\section{Private content and common modular flow}\label{sec:cp}

A private complement \(P_i^{\rm priv}\subseteq A_i\), when one is used,
means \(A_i=W^*(C\cup P_i^{\rm priv})\); no canonical tensor or free
factorization is implied. A state-preserving expectation
\(E_i:A_i\to C\) is a separate regime condition.

\paragraph{Localization of discrepancy.}\label{prop:cpdiscrepancy}
For $\delta=\omega_1|_{A_1\cap A_2}-\omega_2|_{A_1\cap A_2}$,
JS-B directly gives $C\subseteq\ker\delta$, so $\delta(x)\ne0$
implies $x\notin C$.

The kernel need not be an algebra. In \(M_3\), take states with
densities \(\rho=\operatorname{diag}(.5,.3,.2)\),
\(\sigma=\operatorname{diag}(.3,.4,.3)\), and
\(X=\operatorname{diag}(0,1,-1)\). Then
\((\rho-\sigma)(X)=0\) but \((\rho-\sigma)(X^2)=-.2\).
Nor do matching marginals specify a joint-record agreement rate.

\begin{corollary}[Common modular flow]\label{cor:sharedflow}
If faithful normal \(E_i:A_i\to C\) preserve \(\omega_i\), then
\[
 \sigma_t^{\omega_1}|_C=\sigma_t^\psi=\sigma_t^{\omega_2}|_C .
\]
It equals the reference flow used for a Cut if
\(\psi=\omega_0|_C\); equality of flows need not imply equality of states.
For example, on \(C=\mathbb C^2\), the faithful states with weights
\((1/3,2/3)\) and \((2/3,1/3)\) differ but both modular flows are
trivial. No factor condition is assumed here.
\end{corollary}
\begin{proof}
Apply Takesaki's theorem to each expectation and its preserved state.
\end{proof}

\begin{example}[A Type~III shared sector with different source states]
Let \(B=\mathrm{Vir}_1(0,\infty)\), with vacuum state \(\omega_B\);
let \(K,P_1,P_2=M_2\) with normalized trace \(\tau\). Put
\(M=B\bar\otimes K\bar\otimes P_1\bar\otimes P_2\),
\(C=B\otimes1\), and
\[
 A_1=B\bar\otimes K\bar\otimes P_1\otimes1,\qquad
 A_2=B\bar\otimes K\otimes1\bar\otimes P_2 .
\]
With faithful \(\eta_i\) and \(\rho_\pm=(1\pm aZ)/2\), \(0<a<1\),
choose
\(\varpi_1=\omega_B\otimes\rho_+\otimes\eta_1\otimes\tau\) and
\(\varpi_2=\omega_B\otimes\rho_-\otimes\tau\otimes\eta_2\).
They agree on \(C\), differ by \(2a\) on \(Z_K\), and their partial
state expectations onto \(C\) preserve the respective states.
The tracial finite-factor reference makes the reference flow on \(C\)
the vacuum dilation flow. Strong additivity of the \(c=1\) Virasoro net
makes \(\mathrm{Vir}_1(b,\infty)\subset B\), \(b>0\), a standard Cut
on its own vacuum space \cite{FoundationI,PaperI}.
This witnesses the algebra, state and Cut conditions together; source
semigroups and formation of foregrounds have not been specified.
\end{example}

\section{Finite families of independent sources}\label{sec:nsource}

Theorem~\ref{thm:main} extends to any finite number of sources by
requiring one candidate \(C\subseteq\bigcap_iA_i\), one shared seed
algebra in \(C\), JS-A/B for every source, and one Cut inside the
resulting hull. Its proof is the same intersection and restriction
argument. If joint projected records are required, pairwise commuting
normal UCP idempotents give an idempotent finite product; algebraic
range and one common kernel still need to be supplied. Pairwise choices
of overlap sectors and kernels do not establish one coherent joint
process. The same-source theorem in the main text already gives the
largest common admissible algebra without these independent-state
conditions.

\section{Imported interfaces}\label{app:premises}

Foundation~I supplies the hull construction, stationary fixed
algebras, restricted modular groups and conditional Cut inheritance
\cite{FoundationI}. Its projected stationary fixed-sector memory
kernel vanishes; a general projected kernel remains extra dynamical
data. Consciousness~I supplies the half-sided modular inclusion/net
correspondence, and Consciousness~II supplies the bounded pulse and
conditional higher-dimensional identification used in
Section~\ref{sec:sharedsource} \cite{PaperI,PaperII}.
T-DOME's record value, history-selected frame and priced calibration
are different functional interfaces \cite{TDOMEI,TDOMEII,TDOMEIII}.
They do not derive JS-B or preparation permission from a fixed
algebra. The no-foreground-overlap example and the independent-state
witnesses answer different questions and cannot replace one another.
 
\section{Simulation protocol for Section~\ref{sec:model}}\label{app:protocol}

The complete protocol is retained in the supplementary source record
\texttt{FII\_protocol.tex}, with the measured tables.
Its implementation is \texttt{d2\_v3.py}; the selected independent
fault checks use \texttt{d2\_checks.py}. The source map identifies the
fixed scripts and outputs. All reads use legal clipped Bernoulli
instruments; fault and no-fault runs share the world and per-reader
random streams. Calibration uses current reference reads before the
current report, while agreement updates the following epoch.
Communication is quantized to eleven bits per scalar. The scorer's
fault-time isolation is kept separate from deployable detection.

\newpage
\addcontentsline{toc}{section}{References}

\begin{thebibliography}{99}
\raggedright
\sloppy

\bibitem{FoundationI}
S.~Liu, \emph{Foundation I: From Source to Manifestation. Accessible Algebras from a Source, Its Noise and a Readable Seed}, 2026; series record DOI: 10.5281/zenodo.19627786.

\bibitem{FoundationIII}
S.~Liu, \emph{Foundation III: The Perfuming Loop. Held Content, Learned Frames and Physical Return}, 2026; series record DOI: 10.5281/zenodo.19642189.

\bibitem{ILP1998}
M.~Izumi, R.~Longo and S.~Popa,
\emph{A Galois correspondence for compact groups of automorphisms of von Neumann algebras with a generalization to Kac algebras},
J.\ Funct.\ Anal.\ \textbf{155}, 25--63 (1998); arXiv:funct-an/9604004.

\bibitem{BRS2007}
S.~D.~Bartlett, T.~Rudolph and R.~W.~Spekkens,
\emph{Reference frames, superselection rules, and quantum information},
Rev.\ Mod.\ Phys.\ \textbf{79}, 555--609 (2007),
\href{https://doi.org/10.1103/RevModPhys.79.555}{doi:10.1103/RevModPhys.79.555}.

\bibitem{PaperI}
S.~Liu, \emph{Geometry Relative to a Cut: Selection, Generation and Modular Inseparability} (Consciousness~I), Zenodo (2026), DOI: 10.5281/zenodo.19607387.

\bibitem{PaperII}
S.~Liu, \emph{Higher-Dimensional Consistency of Geometry and Modular Flow: Physical Source Identification and a Second-Order Extension}, 2026; series record DOI: 10.5281/zenodo.19607543.

\bibitem{HislopLongo1982}
P.~D.~Hislop and R.~Longo, \emph{Modular structure of the local
algebras associated with the free massless scalar field theory},
Commun.\ Math.\ Phys.\ \textbf{84} (1982) 71--85.

\bibitem{BGL1993}
R.~Brunetti, D.~Guido and R.~Longo, \emph{Modular structure and
duality in conformal quantum field theory},
Commun.\ Math.\ Phys.\ \textbf{156} (1993) 201--219;
\url{https://arxiv.org/abs/funct-an/9302008}.

\bibitem{HAFFA}
S.~Liu, \emph{Physical Access and Correlation Geometry: The Holographic Alaya--Field Framework}, Zenodo (2026), DOI: 10.5281/zenodo.18361706.

\bibitem{HAFFD}
S.~Liu, \emph{Selection, Maintenance and Gravitational Response: Modular Constraints and a Physical Return}, 2026; series record DOI: 10.5281/zenodo.18388881.

\bibitem{TDOMEI}
S.~Liu, \emph{Memory as a Usable Resource: Correlations, Records and Complete Endpoint Accounting}, 2026; series record DOI: 10.5281/zenodo.18574342.

\bibitem{TDOMEII}
S.~Liu, \emph{Reference Frames Selected by History: The Capability and Cost of a Bounded Agent's Self-Frame}, 2026; series record DOI: 10.5281/zenodo.18579703.

\bibitem{TDOMEIII}
S.~Liu, \emph{Keeping a Self-Model Calibrated: Drift Detection, the Price of Keeping Up, and Reliability in a Closed Loop}, 2026; series record DOI: 10.5281/zenodo.18591771.

\bibitem{Nakajima1958}
S.~Nakajima,
\emph{On quantum theory of transport phenomena},
Prog.\ Theor.\ Phys.\ \textbf{20}, 948--959 (1958), DOI:~10.1143/PTP.20.948.

\bibitem{Zwanzig1960}
R.~Zwanzig,
\emph{Ensemble method in the theory of irreversibility},
J.\ Chem.\ Phys.\ \textbf{33}, 1338--1341 (1960), DOI:~10.1063/1.1731409.

\bibitem{Hiai2021}
F.~Hiai, \emph{Concise lectures on selected topics of von Neumann algebras}, arXiv:2004.02383 (2020).

\bibitem{Takesaki1972}
M.~Takesaki,
\emph{Conditional expectations in von Neumann algebras},
J.\ Funct.\ Anal.\ \textbf{9}, 306--321 (1972).

\bibitem{AccardiCecchini1982}
L.~Accardi and C.~Cecchini,
\emph{Conditional expectations in von Neumann algebras and a theorem of
Takesaki},
J.\ Funct.\ Anal.\ \textbf{45}, 245--273 (1982).

\bibitem{Popa1983}
S.~Popa,
\emph{Orthogonal pairs of \(*\)-subalgebras in finite von Neumann
algebras},
J.\ Operator Theory \textbf{9}, 253--268 (1983).

\bibitem{Zurek2009}
W.~H.~Zurek,
\emph{Quantum Darwinism},
Nature Phys.\ \textbf{5}, 181--188 (2009).
\href{https://doi.org/10.1038/nphys1202}{doi:10.1038/nphys1202}.

\bibitem{BrandaoPianiHorodecki2015}
F.~G.~S.~L.~Brand\~ao, M.~Piani, and P.~Horodecki,
\emph{Generic emergence of classical features in quantum Darwinism},
Nat.\ Commun.\ \textbf{6}, 7908 (2015), DOI:~10.1038/ncomms8908.

\bibitem{KorbiczHorodeckiHorodecki2014}
J.~K.~Korbicz, P.~Horodecki, and R.~Horodecki,
\emph{Objectivity in a noisy photonic environment through quantum state information broadcasting},
Phys.\ Rev.\ Lett.\ \textbf{112}, 120402 (2014), DOI:~10.1103/PhysRevLett.112.120402.

\end{thebibliography}
 
\end{document}
